\section{Ablation and Analysis}
\label{sec:ablation}

\subsection{Component ablation}

Table~\ref{tab:ablation} reports the ablation on the same 29 functions and 30
seeds. Each row changes exactly one mechanism of the full method.

\begin{table*}[t]
\centering
\caption{Component ablation on CEC2017, $D=50$, $10^{5}$ true evaluations, 30 seeds.}
\label{tab:ablation}
\begin{tabular}{lcccc}
\toprule
Variant & Dual state & Geometric screening & $M>K$ & Quota \\
\midrule
GWO & -- & -- & -- & -- \\
RASE-GWO & \checkmark & \checkmark & \checkmark & Q-learning \\
fixed $5{:}5$ & \checkmark & \checkmark & \checkmark & fixed \\
fixed $9{:}1$ & \checkmark & \checkmark & \checkmark & fixed \\
fixed $1{:}0$ & \checkmark & \checkmark & \checkmark & fixed \\
random screening & \checkmark & -- & \checkmark & fixed \\
unverified leaders & -- & \checkmark & \checkmark & fixed \\
dense synchronous & \checkmark & \checkmark & -- & -- \\
reduced population & -- & -- & -- & -- \\
\bottomrule
\end{tabular}
\end{table*}

\subsection{Over-sampling factor}

The evaluation ratio $\rho = K/M$ controls how many candidates compete for the
$K$ evaluation slots. Reducing it (larger $M$ at fixed $K$) selects the quota
from a larger candidate pool at no extra evaluation cost. Table~\ref{tab:msweep}
reports the 11-configuration screen used to choose the default.

\subsection{Crowd-level versus leader-level diversity}

The screening rule inherited from sparse PSO/CSO frameworks spends part of the
quota on individuals far from the population centroid. In a leader-driven
algorithm such an individual only matters if it enters the elite triangle, which
rarely happens for an extreme outlier. Replacing that criterion with
minority-leader probing (using the leader farthest from the consensus centre)
allocates the same quota to a second basin instead. The difference between the
two rules is reported together with the fraction of diversity-quota evaluations
that produce an improvement of the best verified fitness.

\subsection{What the Q-learner learns}

Because the state includes the third of the budget consumed, the learned policy
can be read directly. Averaged over runs and problems the agent converges to an
exploitation-dominant policy, briefly opens the minority-leader quota in the
middle third of the budget, and closes it again in the final third. This is
consistent with the fixed-quota study: the consensus-refinement share should be
high throughout, but a temporary opening is what allows a second basin to be
recovered. The learned split is reported per phase.

\subsection{Negative results}

Three findings constrain the claims of this paper and are reported in full.

\textbf{Sparse evaluation without over-sampling.} Holding the number of
generations equal, spending $K=10$ evaluations per round instead of $K=50$ is
clearly harmful on the 12-function development screen, even with the saved
budget converted into five times as many rounds. The saved evaluations do not
buy back the information lost by evaluating a fifth of the candidates, so the
``fewer evaluations at equal accuracy'' claim that holds for sparse-evaluation
PSO/CSO does not transfer to GWO. This is why the paper is framed around the
allocation of a \emph{fixed} evaluation budget rather than its reduction.

\textbf{Crowd-level diversity.} Scoring exploration candidates by distance from
the population centroid converts very few evaluations into improvements; the
minority-leader criterion is used instead.

\textbf{Population reduction.} Simply running GWO with ten wolves, which
matches the number of evaluations per round of the sparse development screen, is
the weakest configuration tested. Over-sampling the candidate pool and reducing
its size are therefore not interchangeable.

\subsection{Behaviour of the learned policy}
\TODO{insert the per-phase learned share table/figure once the formal batch is
analysed, and the convergence curves from \texttt{traces.mat}.}
